\documentclass{beamer}
\usecolortheme{spruce}

\usepackage{heuristica} % Just for cosmetic purposes


% -- Since .sty and .txt are not in same directory, we have to give the
% -- templatefile argument and use ../gwbar
\newlength{\customheight}
\usepackage[
    templatefile="../code/generic_template.txt",
    % -- In the following we explore many different types of signalheight input
    % signalheight="frametitleheight",
    % signalheight="framesubtitleheight",
    % signalheight=1em,
    signalheight=\customheight,
]{../gwbar}
\setbeamertemplate{frametitle}[gwbar]



\begin{document}

\setlength{\customheight}{1em}  % Only relevant for signalheight=\customheight
\begin{frame}[t]{Frame Title}
This is the first slide
\end{frame}

% \setlength{\customheight}{2em}  % Only relevant for signalheight=\customheight
% -- Here we use the following alternative:
\UpdateGWBarOptions[signalheight=2em]
\begin{frame}[t]{Frame Title}{Frame Subtitle}
Slide number 2
\end{frame}
\UpdateGWBarOptions[%
    signalheight=\customheight,
]  % Resetting to length variable


\setlength{\customheight}{4em}  % Only relevant for signalheight=\customheight
\begin{frame}[t]{Frame Title}
Another slide
\end{frame}

\setlength{\customheight}{0em}  % Only relevant for signalheight=\customheight
% -- Height of 0 is bad, no bar can be seen. If your intention was to
% -- draw a flat line, check out the flat_progressbar file to see how
% -- this can be achieved.
\begin{frame}[t]{Frame Title}{Frame Subtitle}
This is the last slide
\end{frame}

\end{document}
